% !PRINTED 739
\weberpage{739}

%<<CONT
so erhalten wir aus den Gleichungen (13)
\[
2 \sum\limits_{1,\, \nu-1}^{i} a_i\, l_{i,k} = \xi_0 L_{0,k} + \xi_1 L_{1,k} + \cdots + \xi_{\nu'-1} L_{\nu'-1,k}.
\]
Daraus ergiebt sich nach (10)
\[
\sum^{i} a_i\, l_{i,k+\nu'} = -\sum^{i} a_i\, l_{i,k},
\]
und dies bedeutet, dass
\[
\delta = \varepsilon_1^{a_1}\, \varepsilon_2^{a_2} \ldots \varepsilon_{\nu-1}^{a_{\nu-1}}
\]
eine Normaleinheit des Körpers $H_m$ ist. Daraus folgt dann weiter,
weil $\mathcal{E}_0, \mathcal{E}_1, \ldots, \mathcal{E}_{\nu'-1}$ ein Fundamentalsystem von Normaleinheiten
 ist, dass die Zahlen $\xi_0, \xi_1, \ldots, \xi_{\nu'-1}$ gerade sein
müssen. Daraus ergiebt sich:

1. \gesperrt{Das System der Gleichungen (18) hat die Congruenzen
}
\weq{19}{\begin{gathered}
\sum\limits_{1,\, \nu-1}^{i} a_i\, M_{k,i} \equiv 0 \pmod{2} \\
k = 0, 1, \ldots, \nu'-1
\end{gathered}}
\gesperrt{zur Folge.}

In den Gleichungen (18), deren Anzahl nur $\nu'-1$ beträgt,
sind aber $\nu-1$ Unbekannte $a_i$ enthalten. Daher giebt es mehrere
Systeme von einander unabhängiger Lösungen, was wir etwas
genauer dahin präcisiren können:

2. \gesperrt{Es giebt $\nu'$ ganzzahlige Lösungen des Systems
(18):}
\weq{20}{
\begin{aligned}
& a_{1,1}, a_{2,1}, \ldots, a_{\nu-1,1} \\
& a_{1,2}, a_{2,2}, \ldots, a_{\nu-1,2} \\
& \ldots\ldots\ldots\ldots \\
& a_{1,\nu'}, a_{2,\nu'}, \ldots, a_{\nu-1,\nu'}
\end{aligned}
}
\gesperrt{von der Art, dass aus der Matrix (20) eine
Determinante von $\nu'$ Reihen gebildet werden
kann, die eine ungerade Zahl ist.}

Denn es giebt zunächst eine Lösung von (18):
\[
a_{1,1}, a_{2,1}, a_{3,1}, \ldots, a_{\nu-1,1},
\]
in der eine ungerade Zahl vorkommt. Nehmen wir an, was
gestattet ist, da hier auf die Anordnung der Indices nichts an-
%CONT>>
